% !Mode:: "TeX:UTF-8"

\chapter{总结与展望(测试)}

\section{总结(测试)}
论文应有结论。论文的结论是最终的、总体的结论，不是正文中各段的小结的简单重复。结论应该观点明确、严谨、完整、准确、精炼。文字必须简明扼要。如果不可能导出应有的结论，也可以没有结论而进行必要的讨论。可以在结论或讨论中提出建议、研究设想、仪器设备改进意见、尚待解决的问题等。不要简单重复罗列实验结果，要认真阐明本人在科研工作中创造性的成果和新见解，在本领域中的地位和作用，新见解的意义。对存在的问题和不足应作出客观的叙述。应严格区分自己的成果与他人（特别是导师的）科研成果的界限。

\section{展望(测试)}
在学术论文中撰写展望时，应首先客观指出研究的局限性，例如数据量或实验条件有限、方法适用范围受限或模型假设不完全符合实际情况；随后提出未来可行的研究方向，包括方法或模型的优化、算法和参数改进、数据或实验扩展，以及在新的问题、场景或领域中的潜在应用；最后强调这些改进和拓展的学术价值或实际意义，使读者看到研究的潜力和发展空间。写作时应逻辑清晰、紧密结合论文内容、用词保持谦虚而积极，并尽量具体可行，避免空泛表述，从而在一段话中完整呈现局限、改进方向与未来价值。